% Draft for the CVPR paper. Replace target counts with frozen release counts
% after media/state QA and source-group split assignment are complete.
\subsection{Dataset Overview}
\label{sec:dataset}

We introduce \textit{EverydayBifurcationBench}, a benchmark for studying
whether video models respond correctly to controlled changes in physical
conditions and predict the resulting interaction outcomes. The benchmark is
organized around 20 everyday rigid-body event types (T01--T20), including
multi-stage cascades, relay interactions, and multi-object events. For each
event, the data specification defines controlled scene configurations across
four backgrounds and five gravity levels. The planned corpus combines
simulations from Genesis and Isaac, with target capacities of 52,800 and
20,160 videos, respectively (72,960 videos in total). These figures describe
the design target; the released counts should be reported from the frozen,
quality-audited manifests.

Each physical trajectory is rendered from three synchronized fixed camera
views. The views share the same simulation state and therefore constitute
multiple observations of one trajectory, rather than independent physical
trials. Depending on the event, controlled variables include gravity,
initial velocity or impulse, object geometry and mass, friction, restitution,
and support or contact conditions. In addition to broad parameter coverage,
the design includes critical-boundary families: configurations immediately
below, at, and above a candidate threshold, with repeated simulations where
specified. Such families make it possible to study whether a model preserves
the change in downstream outcomes induced by a localized intervention.

The benchmark pairs rendered videos with structured supervision. Per-frame
records include object position, orientation, linear and angular velocity,
and contact information; event records describe key interactions, their
order, and terminal outcomes. Configuration metadata records the simulation
conditions and intervention-family membership. This representation supports
separate evaluation of parameter response, event-branch prediction, and
downstream event sequences. The accompanying PhysEdit pipeline runs video
generation experiments, while GravTrace provides one quantitative measure of
gravity inferred from generated or reference video. GravTrace is one
measurement component of the benchmark and is not intended to summarize all
physical mechanisms or event-level correctness.

For the planned joint Genesis--Isaac protocol, the 16 training-pool event
types are split by complete source groups into training, validation, and
in-distribution test partitions, with an approximate 80/10/10 allocation.
T03 and T20 are held out as OOD-1 event types, and T11 and T19 as OOD-2 event
types. These are whole-event transfer evaluations: OOD-1 emphasizes unseen
cascade events, whereas OOD-2 emphasizes unseen relay and multi-object
interactions. Both engines occur in the training regime, so these OOD sets do
not constitute engine-held-out evaluation. Source groups bind together
synchronized views, known background variants, repeats and branches from a
critical family, and known corresponding configurations across engines.
Continuous parameter scans are grouped by parameter intervals rather than by
individually sampled points to reduce leakage between nearby configurations.

The target split capacities are 47,040 training, 5,880 validation, 5,880
in-distribution test, 6,240 OOD-1, and 7,920 OOD-2 videos. Whole-group
integrity takes priority over exact per-event ratios. These capacities are
planning targets, not a claim that all target media have passed release QA or
that the final group-level manifests have been frozen. In the paper, we will
report the final number of videos, trajectories, source groups, and QA-passing
samples from the released manifests, and state the associated data version.
